\section{Security of HQB}\label{sec:hqb-proofs}
%
The objective of a hybrid of this form is that the hybrid remains secure so
long as one of the component algorithms remains secure. In other words, one
algorithm may be found to contain weaknesses without compromising the hybrid.
%
The HQB construction is similar to a generic composition of the classic
DH-based KEM in the nominal group and the IND-CCA KEM. Thus, the reader may
ask why the proof of GHP~\cite{...} does not apply. The basic reason this is
true is because HQB is not a fully generic composition of two IND-CCA KEMs:
the traditional DH-based KEM is not IND-CCA as it would be used in
implementations because of x-coordinate truncation: namely, the shared secret
output from the KEM would actually just be the x-coordinate of the group
element $k_2 \defineequals \nexp(\pk_2, \sk_e)$.
%
\paragraph{Proof-strategy.}
%
Our proofs cover multiple scenarios, including the possibility of
cryptographically relevant quantum computers becoming a reality, and thus
breaking the $\indcca$ security of pre-quantum KEMs, and the possibility of
post-quantum cryptography such as a KEM being classically broken, such as by
an attack, vulnerability, or implementation error.

\medskip \noindent
\textsc{Pre-Quantum Security.} This the case where one of the component
KEMs is assumed to be classically $\indcca$-secure but where the post-quantum
KEM may not offer any $\indcca$ security.
%
We prove in the standard model that HQB is $\indcca$-secure when:
%
\begin{enumerate}
\item The KDF is modeled as a secure $\PRF$;
\item The pre-quantum component is $\sdh$-secure; and
\item The post-quantum KEM is modeled as a possibly broken KEM that provides no security.
\end{enumerate}
%
Our proof also demonstrates the opposite scenario:

\medskip \noindent
\textsc{Post-Quantum Security:} This the case where the component
KEMs is assumed to be $\indcca$-secure against a quantum adversary but the
pre-quantum component group may not offer any security at all.
%
We prove in the standard model that HQB is $\indcca$-secure when:
%
\begin{enumerate}
\item The KDF is modeled as a $\prf$;
\item The pre-quantum component group no longer provides any security; and
\item The post-quantum KEM provides $\INDCCA$ security against a quantum attacker.
\end{enumerate}
%
These results do not rely on the ROM nor the QROM: it is sufficient to assume
the standard model $\PRF$ security property for the employed hash function
$F$.
%
\subsection{Reduction to $\sdh$ security in the ROM}

\begin{theorem} \label{theorem:sdh_kem-coll_imply_hqb}
    %
    Let $\nominalgroup = \nominalgroupexpanded$ be a nominal group and \kem
    be a key encapsulation mechanism, let $H$ be a random oracle with an
    output size of $n$ and let $\adversary{A}$ be an adversary against the
    \cca security of \hqb in Definition \ref{def:hqb} making at most
    $\hashqueries$ queries to the random oracles.
    %
    Then there exists an adversary $\adversary{B}$ such
    that,
    %
    \[ \advantage{\hqb, \adversary{A}}{\cca} \leq 4 \Delta_\nominalgroup +  2 \advantage{\nominalgroup, \adversary{B}}{\sdh} \]
    %
    The run-time of $\adversary{B}$ is roughly the same as of
    $\adversary{A}$.
    %
    $\adversary{B}$ performs at most $2 \hashqueries$ queries to its own
    $\odh$ oracle.
\end{theorem}

\begin{proof}
    We consider the following games:
    %
    % \begin{gamefigure}{Game $G_0 - G_4$.}{fig:sdh_reduction_games}
    %     \begin{games}{2}
    %         \begin{gameenv}{\game $G_0 / \textcolor{\gone}{G_1} / \textcolor{\gtwo}{G_2} / \textcolor{\gthree}{G_3} $}
    %             \State $\ROMtable[\cdot] \assign \bot$
    %             \State $(\sk_1, \pk_1) \assign \kem.\keygen()$
    %             \State $\sk_2 \sample \honestexponent$
    %             \BeginBox[draw=\gone]
    %             \State $\sk_2 \sample \exponent$
    %             \Comment{$G_1 - G_3$}
    %             \EndBox
    %             \State $\pk_2 \assign \nexp(\ngen, \sk_2)$
    %             \State $\pk \assign (\pk_1, \pk_2)$
    %             \State $(k^*_1, c^*_1) \assign \kem.\enc(\pk_1)$
    %             \State $\sk_e \sample \honestexponent$
    %             \BeginBox[draw=\gone]
    %             \State $\sk_e \sample \exponent$
    %             \Comment{$G_1 - G_3$}
    %             \EndBox
    %             \State $c^*_2 \assign \nexp(\ngen, \sk_e)$
    %             \State $k^*_2 \assign \nexp(\pk_2, \sk_e)$
    %             \State $c^* \assign (c^*_1, c^*_2)$
    %             \State $s \assign k^*_1 \concat k^*_2 \concat c^*_1 \concat c^*_2 \concat \pk_1 \concat \pk_2 \concat \text{label}$
    %             \State $\ifthen{\ROMtable[s] = \bot}$
    %             \State \quad $\ROMtable[s] \sample \bits^{n}$
    %             \State $k^* \assign \ROMtable[s]$
    %             \BeginBox[draw=\gthree]
    %             \State $k^* \sample \bits^{n}$
    %             \Comment{$G_3$}
    %             \EndBox
    %             \State $b' \assign \adversary{A}^{\dec(\cdot), H(\cdot)}(\pk, c^*, k^*)$
    %             \State \Return $b'$
    %         \end{gameenv}
    %         \gamebreak
    %         \begin{gameenv}{\oracledec}
    %             \State $\ifthen{c = c^*}$
    %             \State \quad \Return $\bot$
    %             \State $(c_1, c_2) \assign c$
    %             \State $k_1 \assign \kem.\dec(c_1, \sk_1)$
    %             \State $k_2 \assign \nexp(c_2, \sk_2)$
    %             \State $\ifthen{k_1 \compare \bot}$
    %             \State \quad \Return $\bot$
    %             \State $s \assign \hashinput$
    %             \State $\ifthen{\ROMtable[s] = \bot}$
    %             \State \quad $\ROMtable[s] \sample \bits^{n}$
    %             \State $k \assign \ROMtable[s]$
    %             \State \Return $k$
    %             \Statex
    %             \State \underline{\oraclerom{}{m}}
    %             \State $\hashinput \gets m$
    %             \BeginBox[draw=\gtwo]
    %             \LComment{$G_2 - G_3$}\\
    %             $\ifthen{k_2 = k^*_2}$
    %             \State \qquad $\abort_1$
    %             \EndBox
    %             \State $\ifthen{\ROMtable[m] = \bot}$
    %             \State \quad $\ROMtable[m] \sample \bits^{n}$
    %             \State \Return $\ROMtable[m]$
    %         \end{gameenv}
    %     \end{games}
    % \end{gamefigure}

    \begin{itemize}
    \item[\textbf{Game 0}]
        %
        We start with the standard IND-CCA security game for \hqb as defined
        in Figure \ref{fig:sdh_reduction_game_0}.
        %
        In this game, $\Sigma$ is used to model $H(m)$ as lazily sampled random oracle.
        As $G_0$ is the $\cca^0$ game instantiated with \hqb, by definition,
        %
        \[ \prone{\cca^0_{\hqb,\adversary{A}}} = \prone{G^{\adversary{A}}_0}. \]
        %
        \begin{gamefigure}{Game $G_0$.}{fig:sdh_reduction_game_0}
            \begin{games}{2}
                \begin{gameenv}{\game $G_0$}
                    \State $\ROMtable[\cdot] \assign \bot$
                    \State $(\sk_1, \pk_1) \assign \kem.\keygen()$
                    \State $\sk_2 \sample \honestexponent$
                    \State $\pk_2 \assign \nexp(\ngen, \sk_2)$
                    \State $\pk \assign (\pk_1, \pk_2)$
                    \State $(k^*_1, c^*_1) \assign \kem.\enc(\pk_1)$
                    \State $\sk_e \sample \honestexponent$
                    \State $c^*_2 \assign \nexp(\ngen, \sk_e)$
                    \State $k^*_2 \assign \nexp(\pk_2, \sk_e)$
                    \State $c^* \assign (c^*_1, c^*_2)$
                    \State $s \assign k^*_1 \concat k^*_2 \concat c^*_1 \concat c^*_2 \concat \pk_1 \concat \pk_2 \concat \text{label}$
                    \State $\ifthen{\ROMtable[s] = \bot}$
                    \State \quad $\ROMtable[s] \sample \bits^{n}$
                    \State $k^* \assign \ROMtable[s]$
                    \State $b' \assign \adversary{A}^{\dec(\cdot), H(\cdot)}(\pk, c^*, k^*)$
                    \State \Return $b'$
                \end{gameenv}
                \gamebreak
                \begin{gameenv}{\oracledec}
                    \State $\ifthen{c = c^*}$
                    \State \quad \Return $\bot$
                    \State $(c_1, c_2) \assign c$
                    \State $k_1 \assign \kem.\dec(c_1, \sk_1)$
                    \State $k_2 \assign \nexp(c_2, \sk_2)$
                    \State $\ifthen{k_1 \compare \bot}$
                    \State \quad \Return $\bot$
                    \State $s \assign \hashinput$
                    \State $\ifthen{\ROMtable[s] = \bot}$
                    \State \quad $\ROMtable[s] \sample \bits^{n}$
                    \State $k \assign \ROMtable[s]$
                    \State \Return $k$
                    \Statex
                    \State \underline{\oraclerom{}{m}}
                    \State $\hashinput \gets m$
                    \State $\ifthen{\ROMtable[m] = \bot}$
                    \State \quad $\ROMtable[m] \sample \bits^{n}$
                    \State \Return $\ROMtable[m]$
                \end{gameenv}
            \end{games}
        \end{gamefigure}
        %

    \end{itemize}

    \begin{itemize}
    \item[\textbf{Game 1}]
        %
        For $G_1$ we choose the secret keys $\sk_2$ and $\sk_e$ from the set
        of exponents $\exponent$ instead of the set of honest exponents
        $\honestexponent$.
        %
        This allows us to use the $\sdh$ problem to bound the probability of
        the next game hop.
        %
        \begin{gamefigure}{Game $G_1$.}{fig:sdh_reduction_game_1}
        \begin{games}{2}
            \begin{gameenv}{\game $G_1 $}
                \State $\ROMtable[\cdot] \assign \bot$
                \State $(\sk_1, \pk_1) \assign \kem.\keygen()$
                \State $\sk_2 \sample \honestexponent$
                \State \highlightbox{$\sk_2 \sample \exponent$}
                \State $\pk_2 \assign \nexp(\ngen, \sk_2)$
                \State $\pk \assign (\pk_1, \pk_2)$
                \State $(k^*_1, c^*_1) \assign \kem.\enc(\pk_1)$
                \State $\sk_e \sample \honestexponent$
                \State \highlightbox{$\sk_e \sample \exponent$}
                \State $c^*_2 \assign \nexp(\ngen, \sk_e)$
                \State $k^*_2 \assign \nexp(\pk_2, \sk_e)$
                \State $c^* \assign (c^*_1, c^*_2)$
                \State $s \assign k^*_1 \concat k^*_2 \concat c^*_1 \concat c^*_2 \concat \pk_1 \concat \pk_2 \concat \text{label}$
                \State $\ifthen{\ROMtable[s] = \bot}$
                \State \quad $\ROMtable[s] \sample \bits^{n}$
                \State $k^* \assign \ROMtable[s]$
                \State $b' \assign \adversary{A}^{\dec(\cdot), H(\cdot)}(\pk, c^*, k^*)$
                \State \Return $b'$
            \end{gameenv}
            \gamebreak
            \begin{gameenv}{\oracledec}
                \State $\ifthen{c = c^*}$
                \State \quad \Return $\bot$
                \State $(c_1, c_2) \assign c$
                \State $k_1 \assign \kem.\dec(c_1, \sk_1)$
                \State $k_2 \assign \nexp(c_2, \sk_2)$
                \State $\ifthen{k_1 \compare \bot}$
                \State \quad \Return $\bot$
                \State $s \assign \hashinput$
                \State $\ifthen{\ROMtable[s] = \bot}$
                \State \quad $\ROMtable[s] \sample \bits^{n}$
                \State $k \assign \ROMtable[s]$
                \State \Return $k$
                \Statex
                \State \underline{\oraclerom{}{m}}
                \State $\hashinput \gets m$
                \State $\ifthen{\ROMtable[m] = \bot}$
                \State \quad $\ROMtable[m] \sample \bits^{n}$
                \State \Return $\ROMtable[m]$
            \end{gameenv}
        \end{games}
    \end{gamefigure}
        %
    \item[\emph{Claim 1}]
        %
        \[ \gamedifference{0}{1} \leq 2 \Delta_{\nominalgroup}. \]
        %
    \item[\emph{Proof}]
        %
        With this game hop, we change the distribution from which these
        values are chosen.
        %
        $\Delta_{\nominalgroup}$ defines the bound on the probability of an
        adversary to distinguish the original distribution from the new one
        for one element.
        %
        As we are replacing the distribution of two elements, we get $2
        \Delta_{\nominalgroup}$ as an upper bound for $\adversary{A}$ to
        detect this change.
    \end{itemize}

    \begin{itemize}
    \item[\textbf{Game 2}]
        %
        With the changes in $G_2$ we prevent the adversary from querying the
        random oracle containing the value $k^*_2$ that was used to generate
        the challenge key $k^*$.
        %
        This prevents $\adversary{A}$ entirely from obtaining the shared
        challenge key $k^*$ from the random oracle.
        %
        \begin{gamefigure}{Game $G_2$.}{fig:sdh_reduction_game_2}
            \begin{games}{2}
                \begin{gameenv}{\game $G_2$}
                \State $\ROMtable[\cdot] \assign \bot$
                \State $(\sk_1, \pk_1) \assign \kem.\keygen()$
                \State $\sk_2 \sample \honestexponent$
                \State $\sk_2 \sample \exponent$
                \State $\pk_2 \assign \nexp(\ngen, \sk_2)$
                \State $\pk \assign (\pk_1, \pk_2)$
                \State $(k^*_1, c^*_1) \assign \kem.\enc(\pk_1)$
                \State $\sk_e \sample \honestexponent$
                \State $\sk_e \sample \exponent$
                \State $c^*_2 \assign \nexp(\ngen, \sk_e)$
                \State $k^*_2 \assign \nexp(\pk_2, \sk_e)$
                \State $c^* \assign (c^*_1, c^*_2)$
                \State $s \assign k^*_1 \concat k^*_2 \concat c^*_1 \concat c^*_2 \concat \pk_1 \concat \pk_2 \concat \text{label}$
                \State $\ifthen{\ROMtable[s] = \bot}$
                \State \quad $\ROMtable[s] \sample \bits^{n}$
                \State $k^* \assign \ROMtable[s]$
                \State $b' \assign \adversary{A}^{\dec(\cdot), H(\cdot)}(\pk, c^*, k^*)$
                \State \Return $b'$
            \end{gameenv}
            \gamebreak
            \begin{gameenv}{\oracledec}
                \State $\ifthen{c = c^*}$
                \State \quad \Return $\bot$
                \State $(c_1, c_2) \assign c$
                \State $k_1 \assign \kem.\dec(c_1, \sk_1)$
                \State $k_2 \assign \nexp(c_2, \sk_2)$
                \State $\ifthen{k_1 \compare \bot}$
                \State \quad \Return $\bot$
                \State $s \assign \hashinput$
                \State $\ifthen{\ROMtable[s] = \bot}$
                \State \quad $\ROMtable[s] \sample \bits^{n}$
                \State $k \assign \ROMtable[s]$
                \State \Return $k$
                \Statex
                \State \underline{\oraclerom{}{m}}
                \State $\hashinput \gets m$
                \highlightbox{$\ifthen{k_2 = k^*_2}$}
                \State \highlightbox{\qquad $\abort_1$}
                \State $\ifthen{\ROMtable[m] = \bot}$
                \State \quad $\ROMtable[m] \sample \bits^{n}$
                \State \Return $\ROMtable[m]$
            \end{gameenv}
        \end{games}
    \end{gamefigure}
    %
    \item[\emph{Claim 2}]
        %
        \[ \gamedifference{1}{2} \leq Pr[\abort_1] = \advantage{\nominalgroup, \adversary{B}}{\sdh}. \]
        %
    \item[\emph{Proof}]
        %
        To prove this claim, we show the existence of adversary
        $\adversary{B}$ that wins the $\sdh$ game exactly when
        $\adversary{A}$ would query $H$ containing the $k^*_2$ used to
        generate $k^*$.
        %
        Consider adversary $\adversary{B}$ from Figure \ref{fig:sdh} playing
        the $\sdh$ game and simulating $\adversary{A}$'s view on the $\cca$
        game.
        %
        Adversary $\adversary{B}$ uses the challenge group elements $X =
        \nexp(\ngen, x)$ as the static public key and $Y = \nexp(\ngen, y)$
        as the ephemeral public key for the challenge ciphertext.
        %
        Therefore, the shared challenge key $k^*$ is supposed to be $H(
        k^*_1\concat \nexp(Y, x) \concat c^*_1 \concat c^*_2 \concat pk_1
        \concat pk_2\concat\lbl)$.
        %
        $\adversary{B}$ is unable to compute $\nexp(Y, x)$ and therefore sets
        $k^*$ to a uniform random value.
        %
        For $\adversary{A}$ to win the $\cca$ game, it would have to query
        $H$ with $\nexp(Y, x)$.
        %
        $\adversary{B}$ can use the $\odh$ oracle to detect when
        $\adversary{A}$ queries $H$ with the result of $\nexp(Y, x)$, since
        $\odh(Y, k^*_2) = 1$ iff $k^*_2 = \nexp(Y, x)$.

        Therefore, if this check in the $H$ oracle evaluates to $1$ we know
        that the $k_2$ provided has to equal $k^*_2$.
        %
        When this happens, $\adversary{A}$ has provided $\adversary{B}$ with
        the solution to the $\sdh$ game.
        %
        A similar approach can be used to simulate the $\dec$ oracle.
        %
        When calculating the shared key in the $\dec$ oracle, the input for
        the random oracle has to have the form $k_1 \concat \nexp(c_2, x)
        \concat c_1 \concat c_2 \concat pk_1 \concat pk_2 \concat \lbl$ which
        $\adversary{B}$ is not able to compute.
        %
        $\adversary{B}$ can use the same method as before to check whether
        $\adversary{A}$ performs such a query using the $\odh$ oracle.
        %
        We then use the table $\pset{E}$ to ensure consistency between the
        random oracle $H$ and the $\dec$ oracle.
        %
        Therefore,  $\adversary{B}$ simulates  $\adversary{A}$'s view  on the
        $\cca$  game perfectly  and wins  the $\sdh$  game iff  the adversary
        queries $H$  with the result of  $\nexp(Y, x)$, which is  the $k^*_2$
        used to generate $k^*$.
        %
        Note that the adversary $\adversary{B}$ makes at most $2
        \hashqueries$ queries to the $\odh$ oracle.
        %
        \begin{gamefigure}{Adversary $\adversary{B}$.}{fig:sdh}
            \begin{games}{2}
                \begin{gameenv}{\textbf{Adversary} $\adversary{B}^{\odh(\cdot, \cdot)}(X, Y)$}
                    \State $\ROMtable[\cdot] \assign \bot$
                    \State $\pset{E}[\cdot] \assign \bot$
                    \State $(\sk_1, \pk_1) \sample \kem.\keygen()$
                    \State $\pk_2 \assign X$
                    \State $\pk \assign (\pk_1, \pk_2)$
                    \State $(k^*_1, c^*_1) \sample \kem.\enc(\pk_1)$
                    \State $c^*_2 \assign Y$
                    \State $c^* \assign (c^*_1, c^*_2)$
                    \State $k^* \sample \bits^{n}$
                    \State $b' \assign \adversary{A}^{\dec(\cdot), H(\cdot)}(\pk, c^*, k^*)$
                \end{gameenv}
                \gamebreak
                \begin{gameenv}{\oracledec}
                    \State $\ifthen{c = c^*}$
                    \State \quad \Return $\bot$
                    \State $(c_1, c_2) \assign c$
                    \State $k_1 \assign \kem.\dec(c_1, \sk_1)$
                    \State $\ifthen{k_1 \compare \bot}$
                    \State \quad \Return $\bot$
                    \State $\ifthen{\pset{E}[(k_1, \marker, c_1, c_2, \pk_1, \pk_2,\lbl)] = \bot}$
                    \State \quad $\pset{E}[(k_1, \marker, c_1, c_2, \pk_1, \pk_2,\lbl)] \sample \bits^{n}$
                    \State \Return $\pset{E}[(k_1, \marker, c_1, c_2, \pk_1, \pk_2,\lbl)]$
                    \Statex
                    \State \underline{\oraclerom{}{m}}
                    \State $\hashinput \gets m$
                    \State $\ifthen{\odh(Y, k_2) = 1}$ \label{line:found_cdh}
                    \State \qquad output $k_2$
                    \State $\ifthen{\odh(c_2, k_2) = 1}$ \label{line:is_valid_query}
                    \State \qquad $\ifthen{\pset{E}[m] = \bot}$
                    \State \qquad \quad $\pset{E}[m] \sample \bits^{n}$
                    \State \qquad \Return $\pset{E}[m]$
                    \State $\ifthen{\ROMtable[m] = \bot}$
                    \State \quad $\ROMtable[m] \sample \bits^{n}$
                    \State \Return $\ROMtable[m]$
                \end{gameenv}
            \end{games}
        \end{gamefigure}
        %
    \end{itemize}

    \begin{itemize}
    \item[\textbf{Game 3}]
        In $G_3$ we replace the shared challenge key $k^*$ with a uniform
        random key.
        %
        Therefore,
        %
    \item[\emph{Claim 3}]
        %
        \[ \prone{G^{\adversary{A}}_2} = \prone{G^{\adversary{A}}_3}. \]
        %
    \item[\emph{Proof}] To prove that this change is undetectable by
        $\adversary{A}$, we need to argue that $\adversary{A}$'s view when
        interacting with the various oracles is independent of the value of
        $k^*$.
        %
        The abort condition in $H$ prevents the adversary from querying the
        random oracle with the value of $k^*_2$ used to generate the shared
        key $k^*$ from the challenge.
        %
        Since the adversary $\adversary{A}$ cannot query the random oracle
        $H$ with the input used to generate $k^*$, the output of the random
        oracle $H$ is independent of $k^*$.
        %
        The abort condition in \dec also prevents the adversary from querying
        the \dec oracle in a way that would result in $k^*$ being outputted.
        %
        Either $c_1$ or $c_2$ must be different from $c^*_1$ or $c^*_2$.
        %
        If $c_1$ is different, then the resulting key $k_1$ must be different
        from $k^*_1$, otherwise the abort condition would be triggered.
        %
        If $c_2$ is different, then the input to the random oracle must also
        be different, since it is included in its input.
        %
        Therefore, the behavior of the \dec oracle is independent of the
        value of $k^*$.
        %
        Since the behavior of all oracles is independent of $k^*$, the
        adversary cannot detect this game hop.
        %
        \begin{gamefigure}{Game $G_3$.}{fig:sdh_reduction_game_3}
            \begin{games}{2}
                \begin{gameenv}{\game $G_3$}
                    \State $\ROMtable[\cdot] \assign \bot$
                    \State $(\sk_1, \pk_1) \assign \kem.\keygen()$
                    \State $\sk_2 \sample \honestexponent$
                    \State $\sk_2 \sample \exponent$
                    \State $\pk_2 \assign \nexp(\ngen, \sk_2)$
                    \State $\pk \assign (\pk_1, \pk_2)$
                    \State $(k^*_1, c^*_1) \assign \kem.\enc(\pk_1)$
                    \State $\sk_e \sample \honestexponent$
                    \State $\sk_e \sample \exponent$
                    \State $c^*_2 \assign \nexp(\ngen, \sk_e)$
                    \State $k^*_2 \assign \nexp(\pk_2, \sk_e)$
                    \State $c^* \assign (c^*_1, c^*_2)$
                    \State $s \assign k^*_1 \concat k^*_2 \concat c^*_1 \concat c^*_2 \concat \pk_1 \concat \pk_2 \concat \text{label}$
                    \State $\ifthen{\ROMtable[s] = \bot}$
                    \State \quad $\ROMtable[s] \sample \bits^{n}$
                    \State $k^* \assign \ROMtable[s]$
                    \State \highlightbox{$k^* \sample \bits^{n}$}
                    \State $b' \assign \adversary{A}^{\dec(\cdot), H(\cdot)}(\pk, c^*, k^*)$
                    \State \Return $b'$
                \end{gameenv}
                \gamebreak
                \begin{gameenv}{\oracledec}
                    \State $\ifthen{c = c^*}$
                    \State \quad \Return $\bot$
                    \State $(c_1, c_2) \assign c$
                    \State $k_1 \assign \kem.\dec(c_1, \sk_1)$
                    \State $k_2 \assign \nexp(c_2, \sk_2)$
                    \State $\ifthen{k_1 \compare \bot}$
                    \State \quad \Return $\bot$
                    \State $s \assign \hashinput$
                    \State $\ifthen{\ROMtable[s] = \bot}$
                    \State \quad $\ROMtable[s] \sample \bits^{n}$
                    \State $k \assign \ROMtable[s]$
                    \State \Return $k$
                    \Statex
                    \State \underline{\oraclerom{}{m}}
                    \State $\hashinput \gets m$
                    $\ifthen{k_2 = k^*_2}$
                    \State \qquad $\abort_1$
                    \State $\ifthen{\ROMtable[m] = \bot}$
                    \State \quad $\ROMtable[m] \sample \bits^{n}$
                    \State \Return $\ROMtable[m]$
                \end{gameenv}
            \end{games}
        \end{gamefigure}

    \end{itemize}

    \begin{itemize}
    \item[\textbf{Game 4}]
        %
        Finally, we revert the changes introduced in $G_1$ and $G_2$.
        %
        This makes the oracles identical to Game 0, and the only remaining
        change in the main experiment is that $k^*$ is random. Undoing the
        hops accounts for the extra factors in the advantage.
        %
        \begin{gamefigure}{Game $G_4$.}{fig:sdh_reduction_game_4}
            \begin{games}{2}
                \begin{gameenv}{\game $G_4$}
                    \State $\ROMtable[\cdot] \assign \bot$
                    \State $(\sk_1, \pk_1) \assign \kem.\keygen()$
                    \State $\sk_2 \sample \honestexponent$
                    \State $\pk_2 \assign \nexp(\ngen, \sk_2)$
                    \State $\pk \assign (\pk_1, \pk_2)$
                    \State $(k^*_1, c^*_1) \assign \kem.\enc(\pk_1)$
                    \State $\sk_e \sample \honestexponent$
                    \State $c^*_2 \assign \nexp(\ngen, \sk_e)$
                    \State $k^*_2 \assign \nexp(\pk_2, \sk_e)$
                    \State $c^* \assign (c^*_1, c^*_2)$
                    \State $s \assign k^*_1 \concat k^*_2 \concat c^*_1 \concat c^*_2 \concat \pk_1 \concat \pk_2 \concat \text{label}$
                    \State $\ifthen{\ROMtable[s] = \bot}$
                    \State \quad $\ROMtable[s] \sample \bits^{n}$
                    \State $k^* \assign \ROMtable[s]$
                    \State $k^* \sample \bits^{n}$
                    \State $b' \assign \adversary{A}^{\dec(\cdot), H(\cdot)}(\pk, c^*, k^*)$
                    \State \Return $b'$
                \end{gameenv}
                \gamebreak
                \begin{gameenv}{\oracledec}
                    \State $\ifthen{c = c^*}$
                    \State \quad \Return $\bot$
                    \State $(c_1, c_2) \assign c$
                    \State $k_1 \assign \kem.\dec(c_1, \sk_1)$
                    \State $k_2 \assign \nexp(c_2, \sk_2)$
                    \State $\ifthen{k_1 \compare \bot}$
                    \State \quad \Return $\bot$
                    \State $s \assign \hashinput$
                    \State $\ifthen{\ROMtable[s] = \bot}$
                    \State \quad $\ROMtable[s] \sample \bits^{n}$
                    \State $k \assign \ROMtable[s]$
                    \State \Return $k$
                    \Statex
                    \State \underline{\oraclerom{}{m}}
                    \State $\hashinput \gets m$
                    \State $\ifthen{\ROMtable[m] = \bot}$
                    \State \quad $\ROMtable[m] \sample \bits^{n}$
                    \State \Return $\ROMtable[m]$
                \end{gameenv}
            \end{games}
        \end{gamefigure}
        %

    \item[\emph{Claim 4}]
        \[ \gamedifference{3}{4} \leq 2 \Delta_\nominalgroup + \advantage{\nominalgroup, \adversary{B}}{\sdh}. \]

    \item[\emph{Proof}]
        %
        As in $G_1$, we change the distribution from which the two group
        exponents are chosen back to the honest distribution. This incurs $2
        \Delta_{\nominalgroup}$ as an upper bound for $\adversary{A}$ to
        detect this change.
        %
        Undoing the change in $G_2$, we remove the abort in $H$ so
        that a $\odh$ oracle doesn't help anymore, incurring a similar
        advantage difference.
        %
    \end{itemize}

\item
    %
    Now $k^*$ is chosen uniformly at random, which means that $G_4$
    matches the definition of the $\cca^1_{\hqb,\adversary{A}}$ game
    perfectly. Therefore,
    %
    \[ \prone{G^{\adversary{A}}_4} = \prone{\cca^1_{\hqb,\adversary{A}}}. \]

\item
    %
    This concludes the proof of Theorem \ref{theorem:sdh_kem-coll_imply_hqb}.

\end{proof}
